% ======================================================================
% Заключение. Ответственный: A — Плешнев Арсений (собирает выводы от всех).
% Требования (main.tex, строки 1892–1897): основные результаты, расхождения
% и их причины; влияние шага сетки, временного шага и шага оптимизации;
% какие выводы подтверждены аналитически, а какие получены только для
% выбранного численного эксперимента; список внешних источников.
% Образец: report_template.tex, строки 338–351.
% ======================================================================

\todo{A: основные результаты по разделам 1–3 (каждый владелец присылает 2–3 предложения по своей части).}

\todo{A: обнаруженные расхождения и их причины (погрешность квадратур, критерии остановки, машинное округление, критерии обнаружения колебаний).}

\todo{A: влияние шага сетки, временного шага и шага оптимизации на результат.}

\todo{A: отдельно указать, какие выводы подтверждены аналитически, а какие получены только для выбранного численного эксперимента.}

\taskhead{Источник исходных данных}
Условие лабораторной работы \texttt{main.tex} (репозиторий преподавателя),
вариант 3: приложение А — параметры интегралов, приложение Б — параметры модели.

\todo{A: список использованных внешних источников, если были.}
